%%%%%%%%%%%%%%%%%%%%%%%%%%%%%%%%%%%%%%%%%%%%%%%%%%%%%%%%%%%%%%%%%%%%%%%%%%%%%%%%
% DT4TM -- Endabgabe, Version 2
% Stand: 26.09.2026 (Version 2). Lesefassung: project_course_final_v2.md
% Version 1 (project_course_final.tex) bleibt unverändert erhalten.
% Alle Zahlen: report_numbers.py (Kennungen [E1] usw. -> Anhang A).
% Abbildungen: Figures/rom_vs_fem.pdf, Figures/eigen_spectrum.pdf
%              (erzeugt mit: report_numbers.py --figure)
%%%%%%%%%%%%%%%%%%%%%%%%%%%%%%%%%%%%%%%%%%%%%%%%%%%%%%%%%%%%%%%%%%%%%%%%%%%%%%%%

\documentclass[a4paper, 10pt, conference]{ieeeconf}
\IEEEoverridecommandlockouts
\overrideIEEEmargins

\usepackage[utf8]{inputenc}
\usepackage[T1]{fontenc}
\usepackage[ngerman]{babel}
\usepackage{graphicx}
\usepackage{amsmath}
\usepackage[left=1.57cm,right=1.57cm,top=2.54cm,bottom=2.54cm]{geometry}
\usepackage{amssymb}
\usepackage{enumitem}
\usepackage{amsthm}
\usepackage{booktabs}
\usepackage{tikz}
\usetikzlibrary{arrows.meta, positioning, fit, backgrounds, shapes.geometric}
\usepackage{url}
\usepackage{tabularx}
\usepackage[hidelinks]{hyperref}

\theoremstyle{remark}
\newtheorem{remark}{Anmerkung}

\makeatletter
\def\@IEEEprocessthesectionargument#1{%
\@ifmtarg{#1}{%
\@IEEEappendixsavesection*{Anhang \thesectiondis}%
\addcontentsline{toc}{section}{Anhang \thesection}}{%
\@IEEEappendixsavesection*{Anhang \thesectiondis \\* #1}%
\addcontentsline{toc}{section}{Anhang \thesection: #1}}}
\makeatother

% Kasten für Zwischenfazit / Analogie (nur Bordmittel).
\newcommand{\merke}[2]{\par\smallskip\noindent
  \fbox{\parbox{\dimexpr\columnwidth-2\fboxsep-2\fboxrule\relax}{%
  \small\textbf{#1.}~#2}}\par\smallskip}
% Verweis auf eine Zeile von report_numbers.py
\newcommand{\src}[1]{\,\textbf{[#1]}}

\title{\LARGE \bf
Vom Feldmodell zum Echtzeit-Zwilling:\\
Modellordnungsreduktion für den thermischen Digital Twin\\
eines elektrodynamischen Levitators
}
\author{Dang Minh Hoang, Marina Borchert, Mohamed Aziz El Majid
\thanks{Project Course Digital Twin, Summer Term 2026. Supervisors: Hr. Prof. Dr. rer. nat. Dirk Hartmann, Fr. Dr.-Ing. Melina Merkel. Stand: 26.09.2026, Version~2. Zahlen in eckigen Klammern, z.\,B.\ [E3], verweisen auf Anhang~A und die Ausgabe von \texttt{report\_numbers.py}.}%
}

\begin{document}
\maketitle
\thispagestyle{plain}
\pagestyle{plain}

%%%%%%%%%%%%%%%%%%%%%%%%%%%%%%%%%%%%%%%%%%%%%%%%%%%%%%%%%%%%%%%%%%%%%%%%%%%%%%%%
\begin{abstract}
Ein digitaler Zwilling soll die Temperatur von Aluminiumscheibe und Spulen
eines TEAM-28-ähnlichen Levitators laufend aus dem gemessenen Spulenstrom
vorhersagen, auf einem Smartphone. Genaue Simulationsmodelle sind dafür zu
schwerfällig, einfache Faustformeln zu ungenau. Der Bericht zeigt, wie
beides verbunden wird: Die Physik wird als Feldgleichungen formuliert und
mit der Finite-Elemente-Methode (FEM) gelöst ($8700$ komplexe
elektromagnetische und $697$ thermische Unbekannte). Dieses große Modell
wird dann durch Modellordnungsreduktion (MOR) auf acht Zustandsgrößen
verkleinert, deren Zeitschritt $3{,}3\,\mu$s dauert. Die Zulässigkeit
der Reduktion wird nicht angenommen, sondern an drei prüfbaren
Eigenschaften nachgewiesen: Linearität, getrennte Zeitskalen und eine
große Lücke im Eigenwertspektrum ($\tau_1/\tau_2=47$). Gegen die volle FEM
weicht die reduzierte Scheibentemperatur um höchstens $0{,}36\,$K ab. Gegen
den Versuchsstand trifft der Zwilling die kalibrierten stationären
Spulentemperaturen exakt, den Zeitverlauf noch nicht (RMS-Fehler
$11{,}5\,$K); die Ursachen werden benannt. Der Bericht dokumentiert zudem
den Code so, dass seine Funktion ohne Lektüre des Quelltexts
nachvollziehbar ist.
\end{abstract}

%%%%%%%%%%%%%%%%%%%%%%%%%%%%%%%%%%%%%%%%%%%%%%%%%%%%%%%%%%%%%%%%%%%%%%%%%%%%%%%%
\section{Einleitung}
\label{sec:intro}

Der Versuchsstand ist ein elektrodynamischer Levitator nach dem
TEAM-Benchmark~28~\cite{team28}: Zwei mit $50\,$Hz-Wechselstrom betriebene
Spulen erzeugen ein Magnetfeld, das in einer darüber liegenden
Aluminiumscheibe Ströme induziert. Diese werden vom Feld abgestoßen, die
Scheibe schwebt. Dieselben Ströme erzeugen Wärme, in der Scheibe, in den
Spulen und im Eisen.

Genau diese Temperaturen sind im Betrieb schwer zu messen: Ein
Thermoelement an der schwebenden Scheibe stört die Levitation, eine
IR-Kamera misst auf blankem Aluminium unzuverlässig (blankes Metall strahlt
wenig und spiegelt viel), und das Innere der Wicklung ist unzugänglich.
Messbar ist dagegen der Spulenstrom. Daraus ergibt sich die Aufgabe eines
\emph{virtuellen Sensors}: Eine nicht messbare Größe wird aus einer
messbaren über ein physikalisches Modell berechnet. Hartmann und
Van der Auweraer beschreiben dasselbe Muster für die nicht messbare
Rotortemperatur großer Elektromotoren~\cite{hartmann2025golden}. Ein solches
Modell muss drei Anforderungen erfüllen:
\begin{enumerate}[leftmargin=*, nosep]
  \item \textbf{Echtzeitfähigkeit:} mit dem Messtakt ($1\,$Hz) mitlaufen,
        auch auf einem Smartphone im Browser;
  \item \textbf{physikalische Treue:} aus den Grundgleichungen abgeleitet
        sein, damit es auch außerhalb gemessener Fälle verlässlich ist;
  \item \textbf{Kalibrierbarkeit:} die wenigen unsicheren Parameter an
        Messungen nachstellen lassen.
\end{enumerate}
1 und 2 widersprechen sich zunächst: Genaue Modelle sind groß und langsam.
Das Bindeglied ist die \emph{Modellordnungsreduktion}, die Hartmann et al.\
als Schlüsseltechnologie für digitale Zwillinge bezeichnen~\cite{hartmann2018mor}.

\textbf{Aufbau.} Abschnitt~\ref{sec:problem} beschreibt, was modelliert
wird. Abschnitt~\ref{sec:method}, der Kern des Berichts, erklärt, wie und
warum die Gleichungen mit der FEM gelöst und dann reduziert werden.
Abschnitt~\ref{sec:tool} dokumentiert das Berechnungswerkzeug,
Abschnitt~\ref{sec:results} die Ergebnisse im Vergleich mit dem
Versuchsstand, Abschnitt~\ref{sec:outlook} den Ausblick. Anhang~A nennt die
Herkunft jeder Zahl.

%%%%%%%%%%%%%%%%%%%%%%%%%%%%%%%%%%%%%%%%%%%%%%%%%%%%%%%%%%%%%%%%%%%%%%%%%%%%%%%%
\section{Problemstellung}
\label{sec:problem}

\subsection{Physikalisches Modell}
\label{sec:phys}

\textbf{Aufbau.} Abb.~\ref{fig:geometry} zeigt einen halben Querschnitt:
Eisenkern, Innenspule ($1000$ Windungen), Eisenring, Außenspule ($500$
Windungen); darüber schwebt die Aluminiumscheibe (Ø$\,160\,$mm, $3\,$mm dick).
Die Anordnung ist um die senkrechte Achse drehsymmetrisch.

\begin{figure}[t]
\centering
\begin{tikzpicture}[x=0.065cm, y=0.065cm,
    ironfill/.style={fill=gray!45},
    cufill/.style={fill=orange!35},
    alfill/.style={fill=cyan!25},
    lbl/.style={font=\scriptsize, align=center}]
  \draw[dash dot] (0,15) -- (0,-58);
  \node[font=\tiny, right] at (2,16) {$r=0$};
  \draw[-{Stealth[length=4pt]}] (0,-58) -- (109,-58) node[below, font=\tiny] {$r$ [mm]};
  \draw (25.9,-58) -- (25.9,-56.5); \node[font=\tiny] at (25.9,-63) {25{,}9};
  \draw (61.9,-58) -- (61.9,-56.5); \node[font=\tiny] at (61.9,-63) {61{,}9};
  \draw (79.9,-58) -- (79.9,-56.5); \node[font=\tiny] at (79.9,-63) {79{,}9};
  \draw (102.9,-58) -- (102.9,-56.5); \node[font=\tiny] at (102.9,-63) {102{,}9};
  \draw[ironfill] (0,0) rectangle (25.9,-53);
  \node[lbl] at (13,-26.5) {Eisen-\\kern};
  \draw[cufill] (27.9,0) rectangle (61.9,-53);
  \node[lbl] at (44.9,-26.5) {Innenspule\\$N{=}1000$};
  \draw[ironfill] (64.9,0) rectangle (79.9,-53);
  \node[lbl, rotate=90] at (72.4,-26.5) {Eisenring};
  \draw[cufill] (82.9,0) rectangle (102.9,-53);
  \node[lbl] at (92.9,-26.5) {Außen-\\spule\\$N{=}500$};
  \draw[alfill] (0,6) rectangle (80,9);
  \node[lbl] at (40,7.5) {Al-Scheibe, $R{=}80$, $d{=}3$};
  \draw[<->, dashed] (85,0) -- (85,6);
  \node[font=\tiny] at (91,3) {$z_{\mathrm{gap}}$};
  \node[font=\tiny, left] at (0,-26.5) {$53$};
\end{tikzpicture}
\caption{Halber Querschnitt (drehsymmetrisch um $r=0$), Maße in mm aus
\texttt{params.yaml}.}
\label{fig:geometry}
\end{figure}

\textbf{Wirkprinzip.} Der Wechselstrom erzeugt ein Magnetfeld, das sich
$50$-mal pro Sekunde umpolt; Kern und Ring bündeln es. Ein sich änderndes
Feld induziert in jedem Leiter Kreisströme (Wirbelströme), in der Scheibe
und im Eisen. Wirbelströme im Feld erfahren eine Kraft, die die Scheibe
hebt. Zugleich setzen alle Ströme über den elektrischen Widerstand Wärme
frei. Die Wärme verteilt sich durch Leitung und wird an der Oberfläche an
die Luft abgegeben. Gesucht ist die Temperatur überall im Gerät und ihr
zeitlicher Verlauf.

\textbf{Vereinfachungen.} (i) \emph{Drehsymmetrie}: Geometrie und
Anregung ändern sich um die Achse nicht, also genügt die Halbebene $(r,z)$.
(ii) \emph{Umgebung $20\,^\circ$C, Strom $5\,$A} als vereinbarter
Standardfall; beide sind zur Laufzeit änderbar. (iii) \emph{Scheibe in
fester Höhe} $3{,}8\,$mm für die EM-Rechnung; Folgen in
Abschnitt~\ref{sec:causes}.

\textbf{Bekannte und unbekannte Größen.} Tabelle~\ref{tab:inputs} ordnet
die Eingangsgrößen vier Kategorien zu: \emph{gemessen},
\emph{Literatur}, \emph{angenommen} (plausibel, nicht belegt) und
\emph{kalibriert} (aus Messung zurückgerechnet). Sie zeigt die zentrale
Schwierigkeit: Die Wärmeübergänge, die die Endtemperatur bestimmen, sind
geschätzt oder aus einem einzigen Messpunkt zurückgerechnet.

\begin{table}[t]
\centering\footnotesize
\caption{Eingangsgrößen und Herkunft (\texttt{params.yaml}; vollständig in Anhang~A).}
\label{tab:inputs}
\begin{tabularx}{\columnwidth}{@{}X l l@{}}
\toprule
Größe & Wert & Kategorie \\
\midrule
Radien (Abb.~\ref{fig:geometry}), Höhe & $0$--$102{,}9$, $53\,$mm & gemessen \\
Windungen innen/außen & $1000/500$ & Angabe, plausibilisiert \\
Scheibe $R$ / $d$ / Masse & $80\,$mm / $3\,$mm / $159\,$g & gemessen \\
Strom eff. & $5{,}0$ / $7{,}8\,$A & gemessen \\
$\sigma_{\mathrm{Al}}$ / $\sigma_{\mathrm{Cu}}$ & $34$ / $59{,}6\,$MS/m & Literatur \\
$k$, $\rho$, $c_p$ (Al) & $237$, $2700$, $900$ & Literatur \\
$\alpha$ & $3{,}9\cdot10^{-3}\,$K$^{-1}$ & Literatur \\
Eisen $\mu_r$ / $\sigma$ & $1000$ / $1\,$MS/m & \textbf{angenommen} \\
Draht, Füllfaktor & $1{,}2\,$mm, $0{,}6$ & \textbf{angenommen} \\
$h$ Scheibe oben/unten & $10/25\,$W/(m$^2$K) & \textbf{angenommen} \\
Spulenwärme$\to$Luft unten & $0{,}094\,$K/W & \textbf{angenommen} \\
$hA_{\mathrm{innen}}$/$hA_{\mathrm{außen}}$ & $2{,}474/2{,}889\,$W/K & \textbf{kalibriert} \\
IR Spulen bei $7{,}8\,$A & $79/74\,^\circ$C & gemessen \\
\bottomrule
\end{tabularx}
\end{table}

\subsection{Mathematisches Modell}
\label{sec:math}

Die Wirkkette wird Glied für Glied übersetzt; jede Gleichung ist ein
Bilanzsatz.

\textbf{(a) Magnetfeld.} Berechnet wird das magnetische Vektorpotential,
aus dem das Feld durch Ableiten folgt. Wegen der Drehsymmetrie hat es nur
die Umfangskomponente $A_\varphi(r,z)$. Da alles mit derselben Frequenz
sinusförmig schwingt, genügen Amplitude und Phase, zusammengefasst in
einer komplexen Zahl, dem \emph{Phasor}~\cite{biro1989vectorpotential}:
\begin{equation}
  -\nabla\!\cdot(\nu\nabla A_\varphi) + \frac{\nu}{r^2}A_\varphi
  + j\omega\sigma A_\varphi = J_s,\quad \nu=\frac{1}{\mu_r\mu_0}.
  \label{eq:em}
\end{equation}
\emph{Lesehilfe:} Term~1: Ausbreitung des Feldes (Eisen mit großem $\mu_r$
bündelt es). Term~2: geometrischer Preis der Zylinderkoordinaten. Term~3:
Gegenwehr jedes Leiters durch induzierte Wirbelströme. Rechts: Quelle
$J_s=NI/S$ in den Spulen.

\textbf{(b) Wärmequelle.} Die Wirbelstromdichte ist $J=-j\omega\sigma
A_\varphi$; ein Strom setzt pro Volumen $|J|^2/\sigma$ frei, gemittelt über
eine Periode mit Faktor $\tfrac12$:
\begin{equation}
  q(r,z)=\tfrac12\,\sigma\,\omega^2|A_\varphi|^2\;[\mathrm{W/m^3}].
  \label{eq:q}
\end{equation}
Spulen: $P_{\mathrm{Cu}}=\tfrac12I^2R$ mit dem Drahtwiderstand $R$.

\textbf{(c) Temperatur.}
\begin{equation}
  \rho c_p\frac{\partial T}{\partial t}=\nabla\!\cdot(k\nabla T)+q,\quad
  -k\frac{\partial T}{\partial n}=h\,(T-T_\infty).
  \label{eq:heat}
\end{equation}
\emph{Lesehilfe:} Ein Volumenelement erwärmt sich (links), wenn Wärme von
wärmeren Nachbarn zuströmt oder vor Ort entsteht (rechts). Die Oberfläche
gibt umso mehr ab, je wärmer sie gegenüber der Luft ist.

\textbf{(d) Rückkopplung.} Metalle leiten warm schlechter:
\begin{equation}
  \sigma(T)=\frac{\sigma_0}{1+\alpha(T-T_0)} .
  \label{eq:sigmaT}
\end{equation}
$70\,$K Erwärmung senken $\sigma$ um $21\,\%$. Die Wirkung ist
gegenläufig: Scheibenverluste sinken (weniger induzierter Strom),
Spulenverluste steigen (größerer Widerstand bei eingeprägtem Strom).

Die Kette läuft in eine Richtung: \eqref{eq:em} liefert über \eqref{eq:q}
die Quelle für \eqref{eq:heat}; \eqref{eq:heat} wirkt nur über den langsamen
Faktor \eqref{eq:sigmaT} zurück.

%%%%%%%%%%%%%%%%%%%%%%%%%%%%%%%%%%%%%%%%%%%%%%%%%%%%%%%%%%%%%%%%%%%%%%%%%%%%%%%%
\section{Methodik}
\label{sec:method}

Dieser Abschnitt beantwortet zwei Fragen: Wie werden die Gleichungen
gelöst (\ref{sec:fem})? Und warum wird die Lösung reduziert, und woran
erkennt man, dass das zulässig ist (\ref{sec:rom})? Beide Teile folgen
demselben Muster: Argument, Prüfung, Zahl.

\subsection{Finite-Elemente-Methode}
\label{sec:fem}

\textbf{Warum numerisch.} Geschlossene Lösungen von \eqref{eq:em} und
\eqref{eq:heat} gibt es nur für Sonderfälle wie eine Leiterschleife im
unendlichen Raum. Das Gerät vereint fünf Materialbereiche mit Sprüngen der
Permeabilität um den Faktor $1000$ und der Leitfähigkeit von $0$ auf
$34\,$MS/m, in einem begrenzten Gebiet mit Randbedingungen. Man muss das
Gebiet in kleine Stücke zerlegen und die Gleichung stückweise erfüllen.

\textbf{Warum FEM.} Tabelle~\ref{tab:methods} vergleicht die vier gängigen
Verfahren an den Kriterien, die hier zählen. Den Ausschlag geben drei
Punkte: Die FEM behandelt Materialsprünge ohne Zusatzaufwand, weil jedes
Element seinen eigenen Materialwert trägt; EM- und Wärmeproblem lassen sich
mit \emph{demselben} Programmgerüst lösen; und die FEM ist das
Standardwerkzeug der TEAM-Benchmarks, die Ergebnisse sind also mit der
Literatur vergleichbar. Ehrlicherweise: Da die Geometrie aus
Rechteckblöcken besteht, wäre auch ein Differenzenverfahren machbar
gewesen, mit Sonderbehandlung an jeder Materialgrenze.

\begin{table*}[t]
\centering\footnotesize
\caption{Numerische Verfahren im Vergleich (Kriterien dieses Geräts).}
\label{tab:methods}
\begin{tabularx}{\textwidth}{@{}l X X X X@{}}
\toprule
Kriterium & Finite Differenzen & Finite Volumen & Randelemente & \textbf{Finite Elemente} \\
\midrule
Grundidee & Ableitungen durch Differenzen auf einem Gitter & Bilanz über Kontrollvolumen & nur Oberflächen diskretisieren & stückweise einfache Funktionen, Gleichung im gewichteten Mittel \\
Materialsprünge ($\mu_r$: $1\to1000$) & Sonderbehandlung je Grenzfläche & gut & schwierig bei vielen Bereichen & \textbf{natürlich} (Material je Element) \\
Wirbelstromgebiete & möglich & möglich & aufwendig (Volumengebiete) & \textbf{natürlich} \\
Drehsymmetrie ($2\pi r$) & Sonderformeln & möglich & Sonderformeln & \textbf{ein Faktor im Integral} \\
Konvektionsrand & Sonderformeln & gut & möglich & \textbf{natürlich} in der schwachen Form \\
Gleichungssystem & dünn besetzt & dünn besetzt & voll besetzt & dünn besetzt, symmetrisch \\
TEAM-Benchmarks & selten & selten & für Außenraum & \textbf{Standard}~\cite{team28,biro1989vectorpotential} \\
\bottomrule
\end{tabularx}
\end{table*}

\textbf{Schritt~1: Zerlegen.} Die Ebene $(r,z)$ wird in Dreiecke zerlegt
(Abb.~\ref{fig:fem}). In jedem Dreieck ändert sich die Lösung linear, ist
also durch ihre Werte in den drei Ecken (Knoten) bestimmt. Aus einer
unbekannten Funktion werden endlich viele Zahlen. Zu jedem Knoten $i$ gehört
eine „Zeltfunktion“ $v_i$: $1$ am Knoten, linear auf $0$ zu den Nachbarn,
sonst null. Jedes Feld lässt sich als $T\approx\sum_jT_jv_j$ schreiben.

\begin{figure}[t]
\centering
\begin{tikzpicture}[scale=0.95]
  \foreach \i in {0,...,4}{\foreach \j in {0,1,2}{
    \coordinate (p\i\j) at (\i*0.9,\j*0.8);}}
  \foreach \i [evaluate=\i as \ii using int(\i+1)] in {0,...,3}{
    \foreach \j [evaluate=\j as \jj using int(\j+1)] in {0,1}{
      \draw[gray] (p\i\j) -- (p\ii\j) -- (p\ii\jj) -- (p\i\jj) -- cycle;
      \draw[gray] (p\i\j) -- (p\ii\jj);}}
  \fill[orange!35, opacity=0.8] (p11) -- (p10) -- (p20) -- (p31) -- (p32) -- (p22) -- cycle;
  \foreach \i in {0,...,4}{\foreach \j in {0,1,2}{\fill (p\i\j) circle (1.1pt);}}
  \fill[red] (p21) circle (2pt);
  \node[font=\scriptsize, above right] at (p21) {$i$};
  \draw[-Stealth] (-0.3,-0.35) -- (4.0,-0.35) node[right, font=\scriptsize] {$r$};
  \draw[-Stealth] (-0.3,-0.35) -- (-0.3,1.9) node[above, font=\scriptsize] {$z$};
  \begin{scope}[xshift=4.6cm, yshift=0.2cm]
    \draw[-Stealth] (0,0) -- (2.6,0) node[right, font=\scriptsize] {$r$};
    \draw[thick, orange!80!black] (0,0) -- (0.4,0) -- (1.2,1.1) -- (2.0,0) -- (2.4,0);
    \node[font=\scriptsize] at (1.2,1.35) {$v_i$};
    \node[font=\scriptsize, below] at (1.2,0) {$r_i$};
    \draw[dashed] (1.2,0) -- (1.2,1.1);
  \end{scope}
\end{tikzpicture}
\caption{Idee der FEM: unbekannt sind nur Knotenwerte. Die Zeltfunktion
$v_i$ (orange: wo sie ungleich null ist; rechts: Schnitt) koppelt Knoten
$i$ nur mit seinen Nachbarn.}
\label{fig:fem}
\end{figure}

\textbf{Schritt~2: Im Mittel erfüllen (schwache Form).} Mit linearen
Stücken lässt sich \eqref{eq:heat} nicht punktweise erfüllen (zweite
Ableitungen sind null oder unendlich). Man verlangt: Die Gleichung stimmt
im \emph{gewichteten Mittel} um jeden Knoten, gewichtet mit dessen
Zeltfunktion. Physikalisch ist das eine Wärmebilanz um jeden Knoten.
Partielle Integration verschiebt eine Ableitung auf $v_i$, übrig bleiben
wohldefinierte erste Ableitungen:
\begin{multline}
  \underbrace{\int_\Omega k\nabla T\!\cdot\!\nabla v_i\,2\pi r\,dA}_{\text{Leitung}}
  +\underbrace{\oint_\Gamma hTv_i\,2\pi r\,ds}_{\text{Abgabe}}\\
  =\underbrace{\int_\Omega q\,v_i\,2\pi r\,dA}_{\text{Erzeugung}}
  +\oint_\Gamma hT_\infty v_i\,2\pi r\,ds .
  \label{eq:weak}
\end{multline}
Die Konvektionsbedingung erscheint als gewöhnlicher Term; die Achse
$r=0$ braucht im Wärmeproblem keine Bedingung, da $2\pi r$ dort verschwindet.

\textbf{Schritt~3: Gleichungssystem.} Einsetzen von $T=\sum_jT_jv_j$ für
alle Knoten ergibt
\begin{equation}
  \mathbf K\mathbf T=\mathbf f,\qquad
  \text{zeitabhängig:}\ \ \mathbf M\dot{\mathbf T}+\mathbf K\mathbf T=\mathbf f .
  \label{eq:linsys}
\end{equation}
$K_{ij}$ sagt, wie stark Knoten $j$ die Bilanz von Knoten $i$ beeinflusst;
$\mathbf M$ enthält die Wärmekapazitäten. Da sich nur Zeltfunktionen
benachbarter Knoten überlappen, ist $\mathbf K$ fast überall null (dünn
besetzt) und effizient lösbar. Für \eqref{eq:em} läuft alles gleich, mit
komplexer Matrix wegen $j\omega\sigma$. Beide Systeme löst
\texttt{scipy.sparse.linalg.spsolve}~\cite{virtanen2020scipy}; Elementformeln
in Anhang~C.

\textbf{Zwei Entscheidungen, die das Problem klein halten.}
\emph{Phasor statt Zeitschritt:} Eine Feldperiode dauert $20\,$ms, die
Temperatur ändert sich über Minuten, rund $10^4$-mal langsamer
(Anhang~B.1). Die Wärme sieht nur den Periodenmittelwert der Verluste,
genau den liefert \eqref{eq:q}. Eine Zeitbereichssimulation bräuchte für
eine Minute Aufheizen $60\,000$ Schritte (bei $20$ je Periode), nur um
danach zu mitteln. \emph{Grobes Netz in Dickenrichtung:} Ein Wechselfeld
dringt nur bis zur Skintiefe $\delta=\sqrt{2/(\omega\mu\sigma)}$ in einen
Leiter ein; für Aluminium bei $50\,$Hz $12{,}2\,$mm\src{E8}, das Vierfache
der Scheibendicke. Der Wirbelstrom ist über die Dicke nahezu gleichmäßig.
Tabelle~\ref{tab:fem} zeigt die Modellgrößen.

\begin{table}[t]
\centering\footnotesize
\caption{Die beiden FEM-Modelle (Apple M4, Python 3.11, SciPy 1.17; Median aus 5 Läufen).}
\label{tab:fem}
\begin{tabular}{@{}lrr@{}}
\toprule
 & EM-FEM & Wärme-FEM \\
\midrule
Gleichung & \eqref{eq:em}, komplex & \eqref{eq:heat}, reell, SPD \\
Gebiet & $1\times1\,$m & Scheibe \\
Knoten / Dreiecke & $8700$ / $17\,028$\src{E1} & $697$ / $1280$\src{T1} \\
eine Lösung & $0{,}40\,$s\src{E1} & $0{,}03\,$s\src{T1} \\
mit $\mu(B)$-Iteration & $1{,}30\,$s\src{E2} & -- \\
\bottomrule
\end{tabular}
\end{table}

\textbf{Woran man erkennt, dass die Lösung stimmt.} Eine FEM liefert immer
Zahlen, auch bei einem Programmierfehler. Vertrauen entsteht durch
Prüfungen mit bekanntem Sollergebnis:
\begin{enumerate}[leftmargin=*, nosep]
  \item \emph{Energieerhaltung:} erzeugte gleich abgegebene Wärme,
        unabhängig voneinander berechnet: $25{,}81388\,$W zu $25{,}81388\,$W,
        relativ $4\cdot10^{-11}$\src{T2}.
  \item \emph{Gebietsgröße:} Ist das Gebiet groß genug, darf die Art der
        Randbedingung keine Rolle spielen. Dirichlet gegen Neumann: alle
        Verluste $<1\,\%$ verschieden, größte Abweichung Scheibe
        $0{,}77\,\%$\src{E7}.
  \item \emph{Linearität:} doppelter Strom, vierfache Verluste. Linearer
        Löser: $4{,}000$; mit feldabhängigem $\mu(B)$: $4{,}009$ gesamt,
        $3{,}994$ Scheibe\src{E6}, die erwartete leichte Nichtlinearität
        des Eisens.
\end{enumerate}

\merke{Zwischenfazit}{Die FEM macht aus den Feldgleichungen ein großes
lineares Gleichungssystem, eine Unbekannte pro Knoten. Die Lösung ist genau
und geprüft. Jede neue Frage („wie warm bei $6\,$A?“, „in $30\,$s?“)
verlangt aber eine neue Lösung.}

\subsection{Modellordnungsreduktion}
\label{sec:rom}

\textbf{Das Problem.} Ein Zwilling muss \emph{ständig} rechnen (jede
Sekunde ein neuer Messwert), auf \emph{schwacher Hardware} (HTML-Datei im
Handy-Browser, ohne SciPy und ohne Löser für dünn besetzte Systeme), und er
muss \emph{kalibrierbar} sein: Aus einer Messkurve lassen sich keine $697$
Knotenwerte bestimmen, wohl aber wenige deutbare Parameter. Mit voller FEM
kostet jeder Zeitschritt eine EM-Lösung ($1{,}30\,$s\src{E2}) plus eine
Wärmelösung. In 3D wächst die FEM auf $10^5$--$10^7$ Unbekannte und wird
für den Betrieb völlig ungeeignet~\cite{hartmann2025golden}; der 2D-Fall
ist ein überschaubares Lehrbeispiel desselben Problems.

\textbf{Die Grundidee.} Hartmann et al.\ trennen \emph{offline}
(einmal: großes Modell lösen, kleines extrahieren) und \emph{online}
(jeder Zeitschritt: nur das kleine rechnen)~\cite{hartmann2018mor}. Das
funktioniert, weil die Unbekannten nicht unabhängig sind: Heizt man die
Scheibe auf, ändert sich vor allem, \emph{wie stark} sie warm ist, kaum,
\emph{wo}.

\merke{Analogie}{Ein Foto, das nur heller oder dunkler wird: Man speichert
das Bild einmal und merkt sich pro Moment nur einen Helligkeitsfaktor.}

Mathematisch: Das Feld wird aus wenigen festen „Bildern“ (Basisvektoren)
$\mathbf v_k$ mit zeitabhängigen Gewichten $\beta_k$ zusammengesetzt,
$\mathbf T\approx T_\infty+\mathbf V\boldsymbol\beta$, $r\ll N$. Verlangt
man, dass der Fehler senkrecht auf den Bildern steht (dieselbe Idee wie die
schwache Form, nur mit Bildern statt Zeltfunktionen), entsteht ein
$r\times r$-System~\cite{benner2015survey}:
\begin{equation}
  \mathbf V^{\!\top}\!\mathbf M\mathbf V\,\dot{\boldsymbol\beta}
  +\mathbf V^{\!\top}\!\mathbf K\mathbf V\,\boldsymbol\beta
  =\mathbf V^{\!\top}\tilde{\mathbf f}.
  \label{eq:galerkin}
\end{equation}
Die Kunst liegt darin, gute Bilder zu wählen und \emph{nachzuweisen}, dass
wenige genügen.

\textbf{Anzeichen für Reduzierbarkeit.} Nicht jedes Problem ist
reduzierbar, ein turbulentes Strömungsfeld etwa hat Tausende gleich
wichtiger Strukturen. Für dieses Gerät gibt es vier prüfbare Anzeichen
(Tabelle~\ref{tab:signs}), und jedes führt zu einer konkreten Reduktion.

\begin{table*}[t]
\centering\footnotesize
\caption{Anzeichen für Reduzierbarkeit, ihre Prüfung und die resultierende Reduktion.}
\label{tab:signs}
\begin{tabularx}{\textwidth}{@{}l X l l l@{}}
\toprule
Anzeichen & Warum es auf Reduzierbarkeit hinweist & Prüfung & Ergebnis & führt zu \\
\midrule
Linearität im Eingang & Der Strom ändert nur die \emph{Stärke} der Lösung, nie ihre \emph{Form}. & Verluste bei doppeltem Strom & $4{,}000$ (linear), $4{,}009$ mit $\mu(B)$\src{E6} & Reduktion A \\
Getrennte Zeitskalen & Sehr schnelle Vorgänge erscheinen für langsame nur als Mittelwert. & Periode gegen Zeitkonstante & $20\,$ms gegen $204\,$s\src{R5} & Phasor \\
Spektrale Lücke & Klingen alle Muster bis auf eines schnell ab, trägt dieses den langsamen Verlauf. & Eigenwerte der Wärme-FEM & $\tau_1/\tau_2=47$\src{R5} & Reduktion B \\
Kleine Biot-Zahl & Leitet ein Körper innen viel besser, als er abgibt, ist er fast überall gleich warm. & $\mathrm{Bi}=h(d/2)/k$ & $1{,}6\cdot10^{-4}\ll0{,}1$\src{R3} & Reduktion B, C \\
\bottomrule
\end{tabularx}
\end{table*}

\subsubsection{Reduktion A: Linearität der Elektromagnetik}
\emph{Argument:} \eqref{eq:em} ist linear in $A_\varphi$, solange $\mu_r$
und $\sigma$ nicht vom Feld abhängen. Doppelte Quelle $J_s\propto I$ heißt
doppeltes $A_\varphi$, überall im selben Maß; nach \eqref{eq:q} also
vierfaches $q$, ebenfalls überall gleich. Die \emph{Form} der Verlustkarte
ist stromunabhängig. \emph{Folgerung:} Die EM-FEM wird einmal bei
$I_{\mathrm{ref}}=5\,$A gelöst, danach gilt
\begin{equation}
  q(r,z;I,\bar T)=\hat q(r,z)\Big(\frac{I}{I_{\mathrm{ref}}}\Big)^{\!2}
  \frac{\sigma(\bar T)}{\sigma(T_{\mathrm{ref}})} .
  \label{eq:scale}
\end{equation}
Auch der letzte Faktor ist ein Skalar, weil die Scheibe fast überall gleich
warm ist; die Temperatur des vorigen Zeitschritts genügt, da das Feld sofort,
die Temperatur über Minuten reagiert. In MOR-Sprache: parametrische
Reduktion mit einem exakt bekannten Basisvektor $\hat q$, aus $8700$
komplexen Unbekannten wird eine Multiplikation. \emph{Grenze:} Sättigung.
Die Spitzenflussdichte im Eisen beträgt am Betriebspunkt ($5\,$A eff.)
$0{,}96\,$T, $64\,\%$ des angesetzten Sättigungswerts $1{,}5\,$T; am
Kalibrierpunkt $7{,}8\,$A aber $1{,}50\,$T\src{E5}, genau an der Grenze
(Anhang~B.5).

\subsubsection{Reduktion B: eine Mode für die Scheibe}
\emph{Argument:} Aluminium leitet sehr gut, die Scheibe ist dünn, die Luft
nimmt langsam ab: $\mathrm{Bi}=1{,}6\cdot10^{-4}$\src{R3}.
Temperaturunterschiede in der Scheibe gleichen sich viel schneller aus, als
die Scheibe Wärme abgibt; das Feld ist fast gleichmäßig. Die FEM bestätigt:
Übertemperatur überall zwischen $42{,}50$ und $43{,}34\,$K\src{R2}.

\emph{Basisvektor:} Der Code verwendet das stationäre FEM-Feld beim
Referenzstrom, $\mathbf v_1=\Delta\mathbf T_{\mathrm{ref}}$ (ein
\emph{Snapshot}). Mit $r=1$ wird \eqref{eq:galerkin} zu
\begin{equation}
  \tau\dot\beta=\Big(\frac{I}{I_{\mathrm{ref}}}\Big)^{\!2}s(\beta)-\beta,
  \quad T=T_\infty+\beta(t)\,\Delta T_{\mathrm{ref}}(r,z).
  \label{eq:rom1}
\end{equation}
\emph{Lesehilfe:} $\beta=0$ kalt, $\beta=1$ so warm wie im
Referenzfall. Der erste Term ist das Ziel ($\propto I^2$, gedämpft durch den
Leitfähigkeitsfaktor $s$), der zweite die Wärmeabgabe. Nach der Zeit $\tau$
sind $63\,\%$ des Wegs zurückgelegt. Im Code ist $\tau=C/UA$ mit
$C=\rho c_pV=146{,}57\,$J/K und $UA=P_{\mathrm{ref}}/\overline{\Delta
T}=0{,}600\,$W/K, also $\tau=244{,}3\,$s\src{R1} (Anhang~B.3).

\emph{Prüfung 1: Genügt eine Mode?} Der eigentliche Nachweis ist eine
Eigenwertanalyse. Jedes Temperaturfeld der Scheibe zerfällt in
\emph{Eigenformen}, Muster, die beim Abkühlen ihre Form behalten und mit
eigener Zeitkonstante $\tau_i$ abklingen; sie lösen $\mathbf
K\boldsymbol\varphi_i=\lambda_i\mathbf M\boldsymbol\varphi_i$,
$\tau_i=1/\lambda_i$ (Anhang~B.4). Abb.~\ref{fig:spectrum} zeigt: Die
langsamste Form klingt mit $\tau_1=203{,}9\,$s ab, die zweite schon mit
$4{,}38\,$s. Nach etwa $13\,$s ($3\tau_2$) sind alle übrigen unter $5\,\%$
abgeklungen; der minutenlange Verlauf wird allein von der ersten Form
getragen. Diese \emph{spektrale Lücke} $\tau_1/\tau_2=47$ ist der
mathematische Grund für $r=1$.

\begin{figure}[t]
\centering
\includegraphics[width=\columnwidth]{Figures/eigen_spectrum.pdf}
\caption{Zeitkonstanten der sechs langsamsten Eigenformen der Wärme-FEM,
logarithmische Achse\src{R5}.}
\label{fig:spectrum}
\end{figure}

\begin{figure}[t]
\centering
\includegraphics[width=\columnwidth]{Figures/rom_vs_fem.pdf}
\caption{Aufheizen der Scheibe bei $I_{\mathrm{ref}}$. Schwarz: volle FEM
($697$ Unbekannte, implizites Euler, $\Delta t=1\,$s). Blau: eine Mode mit
$\tau_1$. Orange: eine Mode mit $\tau=C/UA$ wie im Code. Unten: größter
Fehler über alle Knoten\src{R6}.}
\label{fig:romfem}
\end{figure}

\emph{Prüfung 2: Wie groß ist der Fehler?} Abb.~\ref{fig:romfem}
vergleicht mit der vollen zeitabhängigen FEM. Mit $\tau_1$ beträgt der
größte Fehler über alle Knoten und Zeiten $0{,}36\,$K, $0{,}8\,\%$ der
Endtemperatur\src{R6}: $697\to1$ Unbekannte kostet praktisch keine
Genauigkeit. Das im Code verwendete $\tau$ liegt $20\,\%$ höher und
erzeugt beim Aufheizen bis $2{,}98\,$K ($6{,}9\,\%$), am Endwert keinen Fehler.

\emph{Warum zwei Zeitkonstanten?} $UA$ wird gegen $20\,^\circ$C gebildet,
die Unterseite gibt ihre Wärme aber an von den Spulen auf
$30{,}0\,^\circ$C vorgewärmte Luft ab\src{T3}. Ein Teil der
Übertemperatur stammt also von der Luft, $UA$ fällt zu klein aus. Die
tatsächliche Wärmeabgabefähigkeit ist $\mathbf 1^{\!\top}\mathbf K\mathbf
1=0{,}719\,$W/K und liefert genau $\tau_1=C/0{,}719$\src{R4}. Der
Zwilling bildet die warme Luft allerdings ohnehin als eigenen, langsam
ansteigenden Anteil ab ($16\,\%$ Luft, $84\,\%$ eigene Wirbelströme\src{R8}),
während die FEM-Referenz sofort warme Luft annimmt. Welche Zeitkonstante
die Wirklichkeit besser trifft, kann nur eine gemessene Aufheizkurve der
Scheibe entscheiden.

\subsubsection{Reduktion C: RC-Netzwerk für Spulen und Eisen}
\emph{Argument:} Gewickeltes Kupfer und kompakter Eisenblock leiten gut;
nach dem Biot-Argument genügt je Körper näherungsweise \emph{eine}
Temperatur. Das ergibt ein thermisches RC-Netz, aufgebaut wie ein
Schaltkreis: Wärmekapazität $\leftrightarrow$ Kondensator, Leitwert
$\leftrightarrow$ Widerstand, Temperatur $\leftrightarrow$ Spannung. Es ist
für elektrische Maschinen etabliert~\cite{wallscheid2021thermal}:
\begin{equation}
  C_i\dot T_i=P_i(I)+\sum_jG_{ij}(T_j-T_i)-hA_i(T_i-T_{\mathrm{air}}).
  \label{eq:rc}
\end{equation}
\emph{Aufbau} (Abb.~\ref{fig:rc}): je Spule eine IR-sichtbare Oberfläche
und ein träges Wicklungsinneres (die Kamera sieht nur die Oberfläche, der
Großteil des Kupfers liegt innen), der Eisenkern und ein gemeinsamer
Luftknoten. $P_i$ stammen aus der EM-FEM und skalieren nach
\eqref{eq:scale}, $C_i$ aus Volumen und Werkstoff; die übrigen Leitwerte
und die Luftknotenwerte sind Schätzungen (Anhang~A).

\begin{figure}[t]
\centering
\begin{tikzpicture}[
  nd/.style={circle, draw, fill=white, align=center, font=\tiny, inner sep=1pt, minimum size=1.25cm},
  src/.style={font=\tiny, align=center},
  lk/.style={font=\tiny, midway, fill=white, inner sep=1pt},
  every edge/.style={draw, thick}]
  \node[nd] (id) at (0,2.6) {Wicklung\\innen\\$854\,$J/K};
  \node[nd] (is) at (0,1.0) {Oberfl.\\innen\\$247\,$J/K};
  \node[nd] (od) at (5.6,2.6) {Wicklung\\außen\\$883\,$J/K};
  \node[nd] (os) at (5.6,1.0) {Oberfl.\\außen\\$255\,$J/K};
  \node[nd] (fe) at (1.9,-0.7) {Eisen\\$1676\,$J/K};
  \node[nd] (air) at (3.9,-0.7) {Luft\\$3000\,$J/K};
  \node[src] (amb) at (3.9,-2.2) {Umgebung $20\,^\circ$C};
  \draw (id) -- node[lk]{$1{,}0$} (is);
  \draw (od) -- node[lk]{$1{,}0$} (os);
  \draw (is) -- node[lk]{$0{,}06$} (fe);
  \draw (is) -- node[lk, pos=0.35]{$2{,}47$} (air);
  \draw (os) -- node[lk]{$2{,}89$} (air);
  \draw (fe) -- node[lk]{$0{,}24$} (air);
  \draw (air) -- node[lk]{$40$} (amb);
  \node[src, left=1mm of is] {$52{,}3\,$W$\,\to$};
  \node[src, right=1mm of os] {$\leftarrow54{,}1\,$W};
  \node[src, left=1mm of fe] {$5{,}9\,$W$\,\to$};
\end{tikzpicture}
\caption{Thermisches RC-Netz (Kreise: Knoten mit Wärmekapazität; Zahlen an
Linien: Leitwerte in W/K; Pfeile: eingespeiste Verluste bei $5\,$A)\src{L1}.
Die Verluste werden im Code massenanteilig auf Oberfläche und Wicklung verteilt.}
\label{fig:rc}
\end{figure}

\emph{Kalibrierung:} $hA_{\mathrm{innen}}$ und $hA_{\mathrm{außen}}$
werden so bestimmt, dass der Zwilling die IR-Messung bei $7{,}8\,$A trifft
(Spulen $79/74\,^\circ$C, Umgebung $29\,^\circ$C\src{N5}). Ein
Nullstellensucher (\texttt{refit\_hA.py}) variiert beide, bis das
\emph{tatsächlich laufende} nichtlineare Modell diese Werte erreicht:
$2{,}474$ und $2{,}889\,$W/K\src{N4}; Kontrolle $79{,}00/74{,}00\,^\circ$C\src{L3}.
Warum eine Handformel früher $7\,$K zu kalt lag, erklärt Anhang~D.

\emph{Grey Box:} Reduktion C ist physikalisch aufgestellt, nicht aus einer
FEM projiziert. Nachteil: keine volle Referenz zum Prüfen. Vorteil: jeder
Parameter ist deutbar und einzeln mess- oder kalibrierbar.

\textbf{Ergebnis der Reduktion.} Aus $8700$ komplexen und $697$ reellen
Unbekannten werden acht Zustände: zwei Gewichte der Scheibe (Wirbelstrom-
und Luftanteil) und sechs Knotentemperaturen. Ein Zeitschritt dauert
$3{,}3\,\mu$s\src{L6}.

\merke{Zwischenfazit}{Die Reduktion ist kein Näherungstrick, sondern nutzt
drei überprüfbare Eigenschaften: Linearität, getrennte Zeitskalen,
spektrale Lücke. Wo eine nachlässt (Sättigung bei hohem Strom), weiß man
genau, wo das reduzierte Modell seine Gültigkeit verliert.}

%%%%%%%%%%%%%%%%%%%%%%%%%%%%%%%%%%%%%%%%%%%%%%%%%%%%%%%%%%%%%%%%%%%%%%%%%%%%%%%%
\section{Berechnungswerkzeug}
\label{sec:tool}

Dieser Abschnitt beschreibt, woher die Daten kommen, in welcher Reihenfolge
welche Module rechnen und wie sie verbunden sind (Abb.~\ref{fig:flow},
Tabelle~\ref{tab:modules}).

\begin{figure*}[t]
\centering
\resizebox{\textwidth}{!}{%
\begin{tikzpicture}[
  node distance=4mm and 6mm,
  box/.style={draw, rounded corners, align=center, font=\footnotesize, fill=white, text width=3.1cm, minimum height=1.05cm},
  db/.style={draw, cylinder, shape border rotate=90, aspect=0.15, align=center, font=\footnotesize, fill=white, text width=2.8cm},
  arr/.style={-Stealth, thick}]
  \node[db] (par) {\texttt{params.yaml}\\\scriptsize Geometrie, Material, Randbed., Kalibrierung, IR-Daten};
  \node[box, right=of par] (cfg) {\texttt{config.py}\\\scriptsize laden, mm$\to$m, $I_{\mathrm{peak}}$};
  \node[box, right=of cfg] (em) {\texttt{em\_solver.py}\\\scriptsize EM-FEM Gl.~(\ref{eq:em},\ref{eq:q}): Verlustkarte, $P$ je Körper, Kraft};
  \node[box, right=of em] (th) {\texttt{thermal\_solver.py}\\\scriptsize Wärme-FEM Gl.~(\ref{eq:weak}), Energiebilanz};
  \node[box, right=of th] (rom) {\texttt{rom.py}\\\scriptsize Red.~A+B: $\Delta T_{\mathrm{ref}}$, $\tau$, $UA$};
  \node[box, below=9mm of em] (lp) {\texttt{lumped\_physics()}\\\scriptsize Red.~C: $C_i$, $G_{ij}$, $hA_i$};
  \node[box, below=9mm of cfg] (fit) {\texttt{refit\_hA.py}\\\scriptsize Kalibrierung $hA$};
  \node[box, below=9mm of rom, fill=cyan!10] (tc) {\texttt{twin\_core.py}\\\scriptsize \texttt{TwinState.step(I,dt)}: Gl.~(\ref{eq:rom1},\ref{eq:rc}) + Schwebehöhe};
  \node[box, below=9mm of th, fill=cyan!10] (js) {JavaScript in \texttt{digital\_twin\_fem.html}\\\scriptsize dieselbe Rechnung};
  \node[db, below=9mm of par] (sen) {Stromsensor\\\scriptsize ACS712 + Arduino, 1\,Hz};
  \node[box, below=7mm of tc] (out) {3D/AR im Browser,\\Desktop-Ansichten};
  \node[box, below=7mm of js, dashed] (xv) {\texttt{xval\_twin.py}\\\scriptsize Python = JavaScript?};
  \draw[arr] (par) -- (cfg); \draw[arr] (cfg) -- (em); \draw[arr] (em) -- (th);
  \draw[arr] (th) -- (rom); \draw[arr] (rom) -- (tc); \draw[arr] (em) -- (lp);
  \draw[arr] (lp) -- (js); \draw[arr] (rom) -- (js);
  \draw[arr] (lp.south) |- ++(0,-0.35) -| ([xshift=-6mm]tc.south);
  \draw[arr, dashed] (fit) -- node[left, font=\tiny]{schreibt $hA$} (par);
  \draw[arr] (sen.east) -- ++(0.6,0) |- node[pos=0.75, below, font=\tiny]{$I_{\mathrm{rms}}(t)$ via \texttt{data\_io.py}} ([yshift=-5mm]tc.west);
  \draw[arr] (tc) -- (out); \draw[arr] (js.east) -- (tc.west |- js.east);
  \draw[dashed] (xv) -- (js); \draw[dashed] (xv) -| (tc);
  \begin{scope}[on background layer]
    \node[fill=gray!12, rounded corners, fit=(cfg)(rom)(lp)(fit), inner sep=5pt,
      label={[font=\footnotesize\bfseries]above:OFFLINE: einmal pro Geometrie ($\approx1{,}4\,$s)}] {};
  \end{scope}
\end{tikzpicture}%
}
\caption{Datenfluss durch den Code. Grau: Offline-Teil (FEM, Reduktion,
Kalibrierung). Blau: Online-Teil, der in jedem Zeitschritt läuft
($3{,}3\,\mu$s), als Python-Referenz und identisch als JavaScript im Browser.}
\label{fig:flow}
\end{figure*}

\textbf{Datenquellen.} (i) \texttt{params.yaml} ist die einzige Quelle aller
Zahlen; kein Modul enthält feste physikalische Konstanten, eine
Nachvermessung ist eine reine Dateiänderung. (ii) Ein Hall-Stromsensor
(ACS712) am Arduino bildet den Effektivwert über exakt $5$ Netzperioden
($100\,$ms) und sendet einmal pro Sekunde eine CSV-Zeile, die
\texttt{data\_io.py} liest. (iii) IR-Messungen (23.06.2026) dienen zur
Kalibrierung und Validierung, nicht als Laufzeiteingang.

\textbf{Offline-Ablauf} (\texttt{python build\_twin\_html\_fem.py},
$\approx1{,}4\,$s): (1)~\texttt{config.py} liest die Parameter. (2)~\texttt{em\_solver.py}
erzeugt ein Dreiecksnetz des $1\times1\,$m-Gebiets (fein am Gerät, grob
außen), weist jedem Dreieck $\nu$, $\sigma$, $J_s$ zu, löst das komplexe
System, passt $\mu$ im Eisen iterativ an $|B|$ an und integriert
\eqref{eq:q} je Körper. (3)~\texttt{thermal\_solver.py} überträgt die
Verlustkarte auf das Scheibennetz, löst \eqref{eq:weak} und prüft die
Energiebilanz. (4)~\texttt{rom.py} speichert $\Delta T_{\mathrm{ref}}$ und
berechnet $C$, $UA$, $\tau$; eine zweite Rechnung ohne vorgewärmte Luft
trennt Wirbelstrom- ($84\,\%$) und Luftanteil ($16\,\%$)\src{R8}.
(5)~\texttt{lumped\_physics()} berechnet das RC-Netz. (6)~Alle Koeffizienten
und das Online-Modell werden als JavaScript in eine HTML-Datei geschrieben,
die ohne eigenen Server läuft (nur three.js wird geladen).

\textbf{Online-Ablauf} (jeder Zeitschritt): (1)~Strom $I$ übernehmen
(Messwert oder Vorgabe). (2)~Verluste nach \eqref{eq:scale} skalieren.
(3)~\eqref{eq:rom1} und \eqref{eq:rc} mit explizitem Euler fortschreiben,
bei $\Delta t>0{,}05\tau$ in Teilschritten. (4)~Schwebehöhe als gedämpfte
Schwingung in geschlossener Form. (5)~$T=T_\infty+\beta\Delta
T_{\mathrm{ref}}$ um die Achse drehen und einfärben.

\begin{table*}[t]
\centering\footnotesize
\caption{Module, Aufgabe, Ein-/Ausgaben (Reihenfolge = Datenfluss).}
\label{tab:modules}
\begin{tabularx}{\textwidth}{@{}l X X X@{}}
\toprule
Datei & Aufgabe & Eingang $\to$ Ausgang & Methode / Bibliothek \\
\midrule
\texttt{config.py} & laden, normieren, Einheiten & YAML $\to$ \texttt{Config} & PyYAML \\
\texttt{em\_solver.py} & EM-FEM, Verluste, Hubkraft, Sättigungs- und Gebietsprüfung & Geometrie, $I$ $\to$ $A_\varphi$, $q$, $P$ je Körper, $F_z$ & P1-FEM komplex, \texttt{spsolve} \\
\texttt{thermal\_solver.py} & Wärme-FEM Scheibe, Energiebilanz & $q(r,z)\to T(r,z)$ & P1-FEM, \texttt{spsolve}, Interpolation \\
\texttt{rom.py} & Reduktion A+B & FEM $\to$ $\Delta T_{\mathrm{ref}}$, $\tau$, $UA$ & Snapshot, Energiebilanz \\
\texttt{twin\_model.py} & Koeffizienten je Scheibenradius & Config, EM, ROM $\to$ Koeffizienten & Zwischenspeicher \\
\texttt{twin\_core.py} & \textbf{Online-Zwilling} \texttt{TwinState.step} & $I(t)\to\beta$, $T_i$, Höhe & expl. Euler; nur NumPy + stdlib \\
\texttt{refit\_hA.py} & Kalibrierung $hA$ durch das laufende Modell & IR-Werte $\to$ \texttt{params.yaml} & \texttt{fsolve} \\
\texttt{data\_io.py} + Arduino & Messstrom einlesen, Zwilling antreiben & seriell/CSV $\to I_{\mathrm{rms}}(t)$ & Effektivwert über 5 Perioden \\
\texttt{build\_twin\_html\_fem.py} & Koeffizienten + Modell in HTML backen & $\to$ \texttt{digital\_twin\_fem.html} & prozedurale 3D-Geometrie, three.js \\
\texttt{xval\_twin.py} & Python- gegen JavaScript-Zwilling & beide $\to$ PASS/FAIL & Playwright, Toleranz $10^{-9}$ \\
\texttt{report\_numbers.py} & alle Zahlen dieses Berichts & Code + Parameter $\to$ Anhang~A & -- \\
\bottomrule
\end{tabularx}
\end{table*}

\textbf{Qualitätssicherung.} Das Online-Modell existiert zweimal (Python
und JavaScript). \texttt{xval\_twin.py} öffnet die HTML-Datei in einem
Browser und vergleicht beide über sieben Stromverläufe (absolute Toleranz
$10^{-9}$) und prüft, ob die eingebackenen Koeffizienten noch zu
\texttt{params.yaml} passen. \texttt{twin\_core.py} enthält sechs
Selbsttests; die FEM-Löser prüfen bei jedem Lauf Energiebilanz und
Gebietsgröße.

%%%%%%%%%%%%%%%%%%%%%%%%%%%%%%%%%%%%%%%%%%%%%%%%%%%%%%%%%%%%%%%%%%%%%%%%%%%%%%%%
\section{Ergebnisse}
\label{sec:results}

\subsection{Verluste und Reduktion}

\begin{table}[t]
\centering\footnotesize
\caption{Verluste aus der EM-FEM (Phasoramplitude $5\,$A, $20\,^\circ$C)\src{E3,\,E4}.}
\label{tab:losses}
\begin{tabular}{@{}lrr@{}}
\toprule
Quelle & $P$ [W] & Anteil \\
\midrule
Wirbelstrom Scheibe & 25{,}81 & 18{,}7\,\% \\
Wirbelstrom Eisen & 5{,}86 & 4{,}2\,\% \\
Ohmsch Innenspule & 52{,}32 & 37{,}9\,\% \\
Ohmsch Außenspule & 54{,}12 & 39{,}2\,\% \\
\midrule
Gesamt & 138{,}11 & 100\,\% \\
\bottomrule
\end{tabular}
\end{table}

\begin{table}[t]
\centering\footnotesize
\caption{Volles gegen reduziertes Modell.}
\label{tab:speed}
\begin{tabular}{@{}lrr@{}}
\toprule
 & FEM (offline) & reduziert (online) \\
\midrule
Unbekannte & $8700$ (kompl.) $+\,697$ & $8$ \\
Aufwand & $1{,}30+0{,}03\,$s\src{E2,T1} & $3{,}3\,\mu$s/Schritt\src{L6} \\
Benötigt & SciPy, dünnbes. Löser & Grundrechenarten \\
Läuft auf & Laptop & auch Smartphone \\
Fehler Scheibe & Referenz & $\le0{,}36\,$K ($\tau_1$)\src{R6} \\
\bottomrule
\end{tabular}
\end{table}

Die Spulen erzeugen $77\,\%$ der Wärme (Tabelle~\ref{tab:losses}), passend
zur Beobachtung, dass sie im Wärmebild die heißesten Bauteile sind. Die
Wattzahlen sind um den Faktor $2$ zu klein, weil der Effektivwert $5\,$A als
Amplitude verwendet wird; für die Temperaturen ist das unschädlich, da
dieser konstante Faktor in der Kalibrierung von $hA$ aufgeht (Anhang~D.3).
Die Kraftrechnung verwendet die echte Amplitude $7{,}07\,$A. Das reduzierte
Modell ist um den Faktor $4\cdot10^5$ schneller als eine vollständige
Neulösung\src{L6} (Tabelle~\ref{tab:speed}).

\subsection{Vergleich mit dem Versuchsstand}

Tabelle~\ref{tab:compare} stellt alle verfügbaren Messwerte der Vorhersage
gegenüber und unterscheidet: \emph{kalibriert} (stimmt per Konstruktion),
\emph{Vorhersage} (echter Test) und \emph{IR unzuverlässig}.

\begin{table}[t]
\centering\footnotesize
\caption{Modell gegen Messung\src{L3,\,L4,\,L5,\,V1}.}
\label{tab:compare}
\begin{tabular}{@{}lrrl@{}}
\toprule
Messgröße & Mess. & Modell & Aussage \\
\midrule
Innensp., stat., $7{,}8\,$A & $79\,^\circ$C & $79{,}00$ & kalibriert \\
Außensp., stat., $7{,}8\,$A & $74\,^\circ$C & $74{,}00$ & kalibriert \\
Außensp., Rampe $140\,$s & $40{,}0$ & $33{,}6$ & Vorhersage \\
Außensp., Rampe $200\,$s & $43{,}0$ & $34{,}7$ & Vorhersage \\
Außensp., Rampe $300\,$s & $48{,}5$ & $36{,}2$ & Vorhersage \\
Außensp., Rampe $360\,$s & $55{,}8$ & $40{,}0$ & Vorhersage \\
Außensp., Rampe $450\,$s & $56{,}0$ & $43{,}6$ & Vorhersage \\
Innensp., Rampe $450\,$s & $60{,}8$ & $44{,}5$ & Vorhersage \\
Eisenkern, stat., $7{,}8\,$A & $45$ & $83{,}7$ & IR unzuverl. \\
Scheibe, stat., $7{,}8\,$A & $37$--$40$ & $107{,}4$ & IR unzuverl. \\
sichtbarer Spalt, $5\,$A & $7$--$8\,$mm & $14{,}7\,$mm & Vorhersage \\
\bottomrule
\end{tabular}

\vspace{2pt}
{\footnotesize Temperaturen in $^\circ$C. RMS-Fehler der Außenspulen-Rampe: $11{,}5\,$K\src{L5}.}
\end{table}

Stationär stimmen die Spulen exakt, aber per Konstruktion; ein
unabhängiger stationärer Punkt fehlt. Im Zeitverlauf heizt das Modell zu
allen sechs Zeitpunkten zu langsam auf. Für Scheibe und Eisen sagt es weit
höhere Temperaturen voraus, als die Kamera zeigt, die dort jedoch
unzuverlässig ist. Die Schwebehöhe wird um etwa $7\,$mm überschätzt.

\subsection{Ursachen der Abweichungen}
\label{sec:causes}

\textbf{(a) Zeitverlauf zu langsam.} Die Endtemperatur hängt nur von den
kalibrierten Wärmeübergängen ab. Wie \emph{schnell} sie erreicht wird,
hängt von den Kapazitäten und ihrer Aufteilung auf Oberfläche und
Wicklung ab. Diese sind aus einer reinen Aufheizmessung nicht eindeutig
bestimmbar: Kleine Oberflächenkapazität mit schwacher Kopplung und große mit
starker Kopplung können dieselbe Aufheizkurve erzeugen. Erst eine
Abkühlkurve trennt beide, weil sie zeigt, wie viel Wärme innen gespeichert
war. Hinzu kommen Messunsicherheiten: Zeitstempel $\pm10\,$s aus einem
mündlichen Protokoll, Umgebung $28{,}6$--$33{,}7\,^\circ$C\src{N6}.

\textbf{(b) Wicklungsinneres unplausibel heiß.} Bei $7{,}8\,$A sagt der
Zwilling $178\,^\circ$C voraus\src{L3}, im Bereich der Grenztemperatur
üblicher Lackdrahtisolation und vermutlich zu hoch. Ursache ist der nur als
Größenordnung geschätzte innere Leitwert $G=1{,}0\,$W/K: Bei kalibrierter
Oberfläche muss das Innere umso heißer werden, je kleiner $G$ ist.

\textbf{(c) Scheibe und Eisen.} Die Kamera war auf die Emissivität $0{,}91$
eingestellt; blankes Aluminium liegt größenordnungsmäßig bei $0{,}05$. Sie
sieht überwiegend gespiegelte Umgebung und zeigt zu tiefe Werte, um einen
unbekannten Betrag. Die Modellwerte beruhen auf geschätzten Wärmeübergängen
und geschätzter Luftvorwärmung. Der Vergleich bestätigt und widerlegt daher
derzeit nichts.

\textbf{(d) Schwebehöhe.} Sättigung ist ausgeschlossen (Kraftänderung
$<0{,}1\,\%$). Wahrscheinlichste Ursache ist das angenommene $\mu_r=1000$,
das die Feldbündelung bestimmt; es wurde bewusst \emph{nicht} auf die
Beobachtung zurückgefittet, weil ein solcher Fit andere Fehler verdecken
würde. Zudem wird die Verlustkarte für $3{,}8\,$mm Höhe berechnet, die
vorhergesagte Gleichgewichtshöhe liegt aber bei $11{,}7\,$mm.

\textbf{(e) Zeitkonstante der Scheibe.} Siehe Abschnitt~\ref{sec:rom}: Die
Code-Definition liegt $20\,\%$ über der Eigenwertanalyse.

%%%%%%%%%%%%%%%%%%%%%%%%%%%%%%%%%%%%%%%%%%%%%%%%%%%%%%%%%%%%%%%%%%%%%%%%%%%%%%%%
\section{Ausblick}
\label{sec:outlook}

\textbf{Messen, bevor weiter modelliert wird.} Fast alle Abweichungen haben
dieselbe Wurzel, fehlende Messdaten für unsichere Parameter:
(1)~eine \emph{Abkühlkurve} der Spulen (Ursachen a, b);
(2)~ein \emph{mattschwarzer Messfleck} bekannter Emissivität auf Scheibe
und Eisen (c, e);
(3)~eine \emph{Messung der Eisenpermeabilität}, etwa eine B-H-Kurve (d);
(4)~ein \emph{unabhängiger stationärer Punkt bei $5\,$A}, damit die
Kalibrierung auch geprüft und nicht nur angepasst ist.

\textbf{Vom Modell zum Zwilling.} Der Stromsensor ist softwareseitig
angebunden. Der nächste Schritt ist, gemessene Temperaturen laufend zur
\emph{Zustandsschätzung} zu nutzen, etwa mit einem Kalman-Filter; das
Modell mit acht Zuständen ist dafür passend dimensioniert.

\textbf{Übertragbarkeit.} Das Vorgehen, eine FEM einmal lösen, die
Reduzierbarkeit nachweisen und online ein kleines Modell rechnen, lässt sich
auf 3D übertragen. Snapshot und Eigenwertanalyse wären dort durch
systematische Verfahren wie Krylov-Unterraum-Methoden zu
ersetzen~\cite{hartmann2025golden,hartmann2018mor}.

\textbf{Lessons Learned.} \emph{Gut:} physikalische Vereinfachungen zuerst;
automatische Prüfungen mit bekanntem Ergebnis haben mehrere Fehler
gefunden; eine einzige Parameterdatei. \emph{Schlecht:} Messdaten zu spät
geplant; Effektiv- und Spitzenwert zweimal verwechselt, bis ein eigener
Variablenname (\texttt{I\_peak}) das unterband; Kalibrierung zunächst über
eine Handformel statt durch das laufende Modell; das reduzierte Modell
wurde nicht von Anfang an gegen seine FEM geprüft, die
$20\,\%$-Abweichung fiel erst bei der Arbeit an diesem Bericht auf.

%%%%%%%%%%%%%%%%%%%%%%%%%%%%%%%%%%%%%%%%%%%%%%%%%%%%%%%%%%%%%%%%%%%%%%%%%%%%%%%%
\appendix

\section{Herkunft aller Zahlen}
\label{app:numbers}

Kennungen im Text verweisen hierher. \texttt{report\_numbers.py} rechnet
alle berechneten Werte neu und gibt sie mit derselben Kennung aus
(\texttt{report\_numbers\_output.txt}). Zeiten: Apple M4, Python 3.11.9,
NumPy 2.4.6, SciPy 1.17.1; absolute Zeiten sind rechnerabhängig.

\begin{table*}[t]
\centering\footnotesize
\caption{Herkunft der Zahlen. Oben: Eingangsdaten (nicht berechnet). Unten: berechnete Werte.}
\label{tab:provenance}
\begin{tabularx}{\textwidth}{@{}l l X l@{}}
\toprule
Kenn. & Größe & Wert & Quelle \\
\midrule
N1 & Strom, Frequenz & $5{,}0\,$A eff., $50\,$Hz, $I_{\mathrm{peak}}=7{,}071\,$A & \texttt{params.yaml: excitation}, \texttt{config.py} \\
N2 & Scheibe & $R=80$, $d=3\,$mm, $\sigma=3{,}4\cdot10^7$, $k=237$, $\rho=2700$, $c_p=900$, $\alpha=3{,}9\cdot10^{-3}$ & \texttt{plate\_material} \\
N3 & Randbed. Scheibe & $h=10/25\,$W/(m$^2$K), $T_{\mathrm{amb}}=20\,^\circ$C, $k_{\mathrm{coil}}=0{,}094\,$K/W (angenommen) & \texttt{thermal\_bc} \\
N4 & RC-Netz & $hA=2{,}4744/2{,}8885$ (kalibriert); Luft $C=3000$, $hA_{\mathrm{far}}=40$ (angenommen) & \texttt{lumped\_thermal} \\
N5 & IR stationär $7{,}8\,$A & Spulen $79/74$, Kern $45$, Scheibe $37$--$40$, Umgebung $29\,^\circ$C & \texttt{validation\_data} \\
N6 & IR Rampe & $40/43/48{,}5/55{,}75/56\,^\circ$C bei $140/200/300/360/450\,$s ($\pm10\,$s) & \texttt{validation\_data} \\
-- & Eisen, Draht & $\mu_r=1000$, $\sigma=10^6$, $B_{\mathrm{sat}}=1{,}5\,$T; Draht $1{,}2\,$mm, Füllf.\ $0{,}6$ (angenommen) & \texttt{iron\_core}, \texttt{coils} \\
\midrule
E1 & EM-Netz, Lösung & $8700$ Knoten, $17\,028$ Dreiecke, $0{,}40\,$s & \texttt{em\_solver: solve\_em} \\
E2 & Verlustpfad & $1{,}30\,$s inkl.\ $\mu(B)$ & \texttt{compute\_losses} \\
E3/E4 & Verluste & $25{,}81/5{,}86/52{,}32/54{,}12$, gesamt $138{,}11\,$W & \texttt{compute\_losses}, \texttt{lumped\_physics} \\
E5 & Flussdichte Eisen & $0{,}48\,$T eff./$0{,}68\,$T Spitze bei $5\,$A Ampl.; $0{,}96\,$T ($5\,$A eff.); $1{,}50\,$T ($7{,}8\,$A eff.) & \texttt{check\_saturation}, Anh.~B.5 \\
E6 & $I^2$-Test & $4{,}000$ linear; $4{,}009$ / $3{,}994$ (Scheibe) mit $\mu(B)$ & \texttt{compute\_losses} bei $I$, $2I$ \\
E7 & Gebietsgröße & Scheibe $0{,}767$, Eisen $0{,}598$, Spulen $0{,}000$, gesamt $0{,}118\,\%$ & \texttt{validate\_domain\_size} \\
E8, R3 & Skintiefe, Biot & $12{,}2\,$mm; $1{,}58\cdot10^{-4}$ & Anh.~B.2 \\
T1/T2/T3 & Wärme-FEM & $697/1280$, $0{,}03\,$s; Bilanz rel.\ $4\cdot10^{-11}$; Luft unten $30{,}01\,^\circ$C & \texttt{thermal\_solver} \\
R1/R2 & ROM & $C=146{,}57\,$J/K, $UA=0{,}5999\,$W/K, $\tau=244{,}3\,$s; $\Delta T$ $42{,}50$--$43{,}34\,$K & \texttt{rom: ThermalROM.build} \\
R4/R5 & Eigenanalyse & $\mathbf 1^{\!\top}\mathbf K\mathbf 1=0{,}7188\,$W/K; $\tau_i=203{,}9/4{,}38/1{,}33/0{,}637/0{,}373/0{,}246\,$s & \texttt{report\_numbers} (\texttt{eigsh}) \\
R6 & ROM gegen FEM & $0{,}359\,$K ($0{,}83\,\%$) mit $\tau_1$; $2{,}976\,$K ($6{,}87\,\%$) mit $C/UA$ bei $212\,$s & \texttt{report\_numbers} \\
R8 & Anteile & Wirbelstrom $0{,}8374$, Luft $0{,}1626$ & \texttt{compute\_eddy\_fraction} \\
L3 & stationär & $7{,}8\,$A/$29\,^\circ$C: $79{,}00/74{,}00$, Eisen $83{,}7$, Luft $35{,}8$, Wicklung $177{,}8/176{,}2$, Scheibe $107{,}4$; $5\,$A/$20\,^\circ$C: $44{,}0/41{,}5$, Eisen $46{,}3$, Scheibe $57{,}5\,^\circ$C & \texttt{TwinState} ($2000\times300\,$s) \\
L4/L5 & Rampe & Tab.~\ref{tab:compare}; RMS $11{,}52\,$K & \texttt{TwinState} ($\Delta t=1\,$s) \\
L6 & Online-Schritt & $\approx3{,}3\,\mu$s (Wiederholungen $3{,}25$--$3{,}28$); Beschleunigung $4{,}1\cdot10^5$ & \texttt{TwinState.step} \\
V1 & Schwebehöhe & $z_{\mathrm{eq}}=11{,}7\,$mm (Unterkante), sichtbar $14{,}7\,$mm; $F_g=1{,}598\,$N & \texttt{run\_rig\_validation} \\
\bottomrule
\end{tabularx}
\end{table*}

\section{Herleitungen}
\label{app:derive}

\textbf{B.1 Zeitskalen.} Periode $1/50\,$Hz $=20\,$ms; $\tau_1=204\,$s;
Verhältnis $\approx10^4$.

\textbf{B.2 Skintiefe, Biot-Zahl.}
$\delta=\sqrt{2/(2\pi\cdot50\cdot4\pi\cdot10^{-7}\cdot3{,}4\cdot10^7)}=12{,}2\,$mm,
$d/\delta=0{,}25$. $\mathrm{Bi}=h(d/2)/k=25\cdot0{,}0015/237=1{,}58\cdot10^{-4}$
(größerer der beiden Wärmeübergänge als ungünstigster Fall).

\textbf{B.3 $\tau=C/UA$.} $C=\rho c_p\pi R^2d=2700\cdot900\cdot\pi\cdot0{,}08^2\cdot0{,}003=146{,}57\,$J/K.
Bei $P_{\mathrm{ref}}=25{,}81\,$W ist $\overline{\Delta T}=43{,}03\,$K, also
$UA=0{,}600\,$W/K und $\tau=244{,}3\,$s. Das folgt aus der Energiebilanz
eines einzelnen Körpers $C\dot{\overline T}=P-UA\,\overline{\Delta T}$.

\textbf{B.4 Eigenwertanalyse.} Ohne Quelle gilt $\mathbf M\dot{\mathbf
T}=-\mathbf K\mathbf T$; der Ansatz $\mathbf T=\boldsymbol\varphi
e^{-t/\tau}$ führt auf $\mathbf K\boldsymbol\varphi=\lambda\mathbf
M\boldsymbol\varphi$, $\tau=1/\lambda$. $\mathbf K$, $\mathbf M$ sind
symmetrisch positiv definit, alle $\lambda_i$ reell positiv; jede Lösung
ist $\mathbf T_{\mathrm{ss}}+\sum_ic_i\boldsymbol\varphi_ie^{-t/\tau_i}$.
Abschneiden nach $i=1$ lässt Terme weg, die nach $3\tau_2\approx13\,$s
unter $5\,\%$ liegen. Numerisch: $\mathbf M$ diagonal („lumped mass“),
Symmetrisierung mit $\mathbf M^{-1/2}$, \texttt{scipy.sparse.linalg.eigsh}
mit Shift-Invert um $0$. Für eine fast konstante Mode liefert der
Rayleigh-Quotient $\lambda_1\approx\mathbf 1^{\!\top}\mathbf K\mathbf
1/\mathbf 1^{\!\top}\mathbf M\mathbf 1=0{,}7188/146{,}57$, also
$\tau_1\approx203{,}9\,$s, in Übereinstimmung mit dem Eigenlöser. Der
Unterschied zu B.3: $\mathbf 1^{\!\top}\mathbf K\mathbf 1$ ist die gesamte
Wärmeabgabe je Kelvin; $UA$ wird gegen $20\,^\circ$C gebildet, obwohl die
Unterseite $30\,^\circ$C warme Luft sieht.

\textbf{B.5 Flussdichte.} $B\propto I$ bei konstantem $\mu_r$. Löser bei
Phasoramplitude $5\,$A: $0{,}480\,$T eff., Spitze $0{,}679\,$T. Bei $5\,$A
eff.\ (Amplitude $7{,}07\,$A): $0{,}679\cdot7{,}07/5=0{,}96\,$T. Bei
$7{,}8\,$A eff.\ ($11{,}03\,$A): $1{,}50\,$T. Die lineare Skalierung
überschätzt leicht, da $\mu_r$ nahe der Sättigung sinkt.

\textbf{B.6 Beschleunigung.} $(1{,}30+0{,}03)\,$s$\,/\,3{,}25\,\mu$s
$\approx4{,}1\cdot10^5$.

\section{Details der FEM}
\label{app:fem}

\textbf{Elemente.} P1-Dreiecke; mit konstanten Formfunktionsgradienten
$b_k$, $c_k$, Fläche $A_e$ und Schwerpunktradius $r_c$ ist $\mathbf
K_e=2\pi k(\mathbf b\mathbf b^\top+\mathbf c\mathbf c^\top)A_er_c$.
\textbf{Konvektionsrand} (Kante der Länge $L$, Radien $r_a$, $r_b$):
$\mathbf K_{\mathrm{Kante}}=2\pi h\frac{L}{12}\begin{pmatrix}3r_a+r_b&r_a+r_b\\r_a+r_b&r_a+3r_b\end{pmatrix}$.
\textbf{Netze.} EM: graduiertes Rechtecknetz in Dreiecke geteilt, fein am
Gerät, grob bis $r=0{,}5\,$m, $|z|=0{,}5\,$m; $A_\varphi=0$ auf Achse und
Außenrand (Vergleich mit Neumann). Wärme: Scheibe mit $40\times16$
Teilungen. \textbf{Löser.} Direkte LU (\texttt{spsolve}); EM komplex
assembliert, ohne Aufspaltung in Real- und Imaginärteil. $\mu$ im Eisen
iterativ an $|B|$ angepasst bis $<2\,\%$ Änderung (2 Iterationen).
\textbf{Kraft.} Periodengemitteltes Lorentzkraft-Integral über die
Wirbelströme der Leiter, mit der echten Amplitude $\sqrt2I_{\mathrm{rms}}$.

\section{Kalibrierung und Messdaten}
\label{app:calib}

\textbf{D.1 Messungen.} Zwei IR-Sitzungen am 23.06.2026 (HIKMICRO):
stationär bei $7{,}8\,$A und $29\,^\circ$C; Rampe $5\,$A bis $300\,$s, dann
$7{,}75\,$A bis $450\,$s. Verlässlich sind nur die dunkel lackierten Spulen.

\textbf{D.2 Kalibrieren durch das laufende Modell.} Das Online-Modell
verwendet $hA_{\mathrm{eff}}=hA\,(\Delta T/\Delta T_{\mathrm{cal}})^{0{,}25}$
(freie Konvektion). Eine frühere Kalibrierung über die lineare Bilanz
$T=T_\infty+P/hA$ ergab im nichtlinearen Modell bei $7{,}8\,$A
$72{,}15/67{,}78$ statt $79/74\,^\circ$C, etwa $7\,$K zu kalt
(\texttt{docs/physics.md}, Abschn.~11d). Seitdem: Nullstellensuche durch das
laufende Modell (\texttt{refit\_hA.py}).

\textbf{D.3 Stromkonvention.} Die Verlustkette nutzt $5\,$A (Effektivwert)
als Amplitude, alle Verluste sind um $(\sqrt2)^2=2$ zu klein. Da $hA$ gegen
gemessene Temperaturen kalibriert wird, fällt $hA$ im selben Verhältnis
kleiner aus; die Temperaturen bleiben korrekt. Die Kraft lässt sich nicht so
kalibrieren, die Kraftkette nutzt $I_{\mathrm{peak}}$.

\section{Reproduktion}
\label{app:repro}

Im Repository-Stamm:\\
\texttt{.venv/bin/python projectseminar\_DT4TM\_report/report\_numbers.py [-{}-figure]}\\
\texttt{python refit\_hA.py}\quad\texttt{python xval\_twin.py}\quad\texttt{python twin\_core.py}

\bibliographystyle{IEEEtran}
\bibliography{dt4tm_references}

\end{document}
